% ---
% id: meat-core/prop:retention
% kind: prop
% title: "Unsoundness of dropping a constraint-active sector"
% summary: >-
%   A constraint-active sector adds no entropy deficit yet shapes the
%   statistics: in a two-sector example, forcing it to zero raises the value
%   from 0 to 1 bit per round. Dropping such a sector from the variable list
%   can therefore overstate the key.
% prominence: standard
% status: implemented
% uses: [prop:trivial-alice-system, def:branch, def:admissible-declaration]
% assumes: []
% refs: []
% provenance: newly-derived
% implements: certiqs_key.meat_core.declaration.sectors:build_sector_report
% lean:
% ---

A constraint-active sector contributes nothing to any $K_{\hat c}$, which
might suggest deleting it along with its variable. Its weight, however,
still enters the statistics, and an adversary may place it there.

\begin{proposition}
\label{prop:retention}
There exists a declaration satisfying Definition~\ref{def:admissible-declaration} and a
constraint-active sector $s_0$ such that forcing $\omega_{s_0}=0$ strictly
\emph{increases} $\hlow$. Dropping a constraint-active sector from the
variable list is therefore unsound: it can overstate the extractable key.
\end{proposition}

\begin{proof}
Take $d_A=1$ ($\tau_A=[1]$, equivalently \tref{untrusted-source mode} by
Proposition~\ref{prop:trivial-alice-system}; single Alice outcome with
$M=1$, announcements
trivial, one announcement class, $\varsigma\equiv\mathrm{pass}$,
$\phi\equiv\hat c_0$, $\Chatset=\Chkey=\{\hat c_0\}$,
$\acSet=\{\delta_{\hat c_0}\}$), $\Zval=\{0,1\}$, and
$\Qsys=\Qsys_1\oplus\Qsys_2$ with both sectors one-dimensional. Declare
$\Vset=\{0,1\}$, $r(v)=v$, and
\begin{equation*}
  F^{(1)}_0=F^{(1)}_1=\tfrac12,
  \qquad
  F^{(2)}_0=1,\quad F^{(2)}_1=0 .
\end{equation*}
Sector 2 is key-inert (its \tref{effective outcome set} is $\{v=0\}$, a
singleton), hence constraint-active; sector 1 is key-active. A feasible
state is $\omega=(w_1,w_2)$, $w_i\ge0$, $w_1+w_2=1$, and
$P_{\hat c_0}=1$. A direct computation of \eqref{eq:GTheta} for sector~1
gives, for weight $w_1$, the rank-one operator
$\cG_{\hat c_0,1}(w_1)=w_1\dyad{\eta}$ with
$\ket{\eta}=\tfrac1{\sqrt2}\sum_{v}\ket{v}_R\otimes\ket{x_0,v}_{XY}
\otimes\ket{\kappa_0}$, so
$\cG^{1/\alpha}=w_1^{1/\alpha}\dyad{\eta}$,
$\cZ(\dyad\eta)=\tfrac12\sum_v\dyad{\xi_v}$ for orthonormal
$\ket{\xi_v}$, and
$T=w_1\,\tr[(\tfrac12\sum_v\dyad{\xi_v})^{\alpha}]=w_1\,2^{1-\alpha}$;
hence $f_{\hat c_0,1}(w_1)=w_1(2^{1-\alpha}-1)$ and
$f_{\hat c_0,2}\equiv0$, so
$K=1+w_1(2^{1-\alpha}-1)$ and
$O_{\alpha}=\frac{1}{\alpha-1}\log_2\frac{1}{K}$. The infimum over
$w_1\in[0,1]$ is attained at $w_1=0$ with value $0$; forcing
$\omega_{s_0}=\omega_2=0$ (i.e.\ $w_1=1$) gives the strictly larger value
$\frac{1}{\alpha-1}\log_2 2^{\alpha-1}=1$. The value $0$ is also the
correct one: an adversary that always populates sector~2 makes the raw key
deterministic. Hence the drop overstates the key by one bit per round.
\end{proof}
